\documentclass[11pt]{article}
\usepackage[a4paper,margin=0.9in]{geometry}
\usepackage{booktabs}
\usepackage{amsmath}
\usepackage{amssymb}
\usepackage{graphicx}
\usepackage{microtype}
\sloppy
\title{Comparative Study of Layer-Reduction Approaches}
\author{Sumit Jha}
\date{October 2026}

\begin{document}
\maketitle

\begin{abstract}
We implemented fifteen layer-reduction ideas for GPT-2 (124M) and
SmolLM-135M and evaluated each on wikitext-2 subsets with a fixed PPL
protocol. Four produced a usable signal: the fitted-ternary TerGRASP
bridge (01), the ShortGPT BI baseline (06), SVD low-rank truncation (07),
and the quantized no-op bypass (15). The rest failed by orders of
magnitude (02, 08, 14), stayed at the baseline by construction (05, 11),
or depended on a training budget we could not afford on one MPS laptop
(03, 09, 10, 12, 13). This file walks through all fifteen, then gives the
maths of each of the four working ones and the argument for taking 01 as
the headline.
\end{abstract}

\section{Setup used for every approach}
Models: GPT-2 small and SmolLM-135M, loaded with \texttt{use\_cache=False}.
Data: Salesforce/wikitext, config wikitext-2-raw-v1 unless noted; train
capped at 300k chars, validation/test at 100k chars. Evaluation:
\[
\mathrm{PPL}=\exp\!\Big(\tfrac{1}{N}\sum_{i=1}^{N} -\log
p_\theta(x_i \mid x_{<i})\Big)
\]
computed over non-overlapping 256-token windows, first 3000 tokens.
Every number in this document is the PPL after the approach is applied to
the lowest-BI layers (see Section~\ref{sec:bi}).

\section{BI ranking, common to most approaches} \label{sec:bi}
For decoder block $l$ we record
\[
\mathrm{BI}_l \;=\; 1 - \mathbb{E}_{t}\big[\cos(h^{(l)}_t, h^{(l+1)}_t)\big]
\]
averaged over the batch and sequence. Low $\mathrm{BI}_l$ means the block
is near the identity $h^{(l+1)}\approx h^{(l)}$, so it is the natural
first candidate for deletion, merging, or aggressive quantization. GPT-2's
five most redundant blocks are 4, 1, 2, 3, 5; SmolLM's are 2, 27, 5, 4, 3.

\section{What we ran and how each came out}

\subsection{01 - Fitted-ternary TerGRASP (WORKS, headline)}
\textbf{Idea.} For each redundant block $L$, collect $X=h_{in}, Y=h_{out}$
on $\sim$4k calibration chars, fit
\[
\min_{U,V}\big\| X\,U\,V - \Delta \big\|_F^2,\quad \Delta = Y-X,
\]
then AbsMean-ternarize $U,V$:
\[
\gamma=\tfrac{1}{mn}\|W\|_1,\quad
W_q=\gamma\,\mathrm{clip}\!\big(\mathrm{round}(W/\gamma),-1,1\big),
\quad W_q\in\{-\gamma,0,\gamma\}.
\]
The block is replaced by $g(x)=x+x\,U_q\,V_q$.
\textbf{Why it works.} At a low-BI position $\Delta$ is small and lives on
a low-dimensional manifold of directions, so a rank-32 factor recovers
most of it; the ternary noise adds $\mathcal{O}(\gamma^2\|x\|)$.
\textbf{Result.} GPT-2 $k{=}1$: 45.11 vs 45.76 hard removal; SmolLM
$k{=}1$: 21.71 vs 19.93. Tracks hard removal for $k\ge 2$.

\subsection{06 - ShortGPT BI baseline (WORKS, the control)}
\textbf{Idea.} Remove the $k$ lowest-BI blocks; the control removes $k$
random ones.
\textbf{Result.} k=1 is free either way (GPT-2 45.76 vs 44.42 random);
from k=2 on targeted removal dominates: GPT-2 k=2 47.95 vs 139.09,
k=3 121.81 vs 4181.21. SmolLM k=4 34.20 vs 57.37 random.
\textbf{Role.} The falsification control: every other approach must beat
or match its BI-targeted curve.

\subsection{07 - SVD low-rank truncation (WORKS on SmolLM)}
\textbf{Idea.} For each Linear $W$ in the two most redundant blocks,
$W\approx U_r S_r V_r^\top$ with rank $r$.
\textbf{Maths.} By Eckart-Young the reconstruction error is the tail
$\sum_{i>r}\sigma_i^2$ of the singular values; it is small when $W$ is
numerically low-rank, which empirically coincides with low-BI blocks.
\textbf{Result.} SmolLM r=256: 19.93 (baseline), r=128: 22.00, r=64:
24.22. GPT-2 rows are a dtype-guard artifact and were archived for a clean
fp32 rerun.

\subsection{15 - Quantized no-op bypass (WORKS, supporting rule)}
\textbf{Idea.} Quantize every weight inside the three lowest-BI blocks to
$b\in\{8,4,2\}$ bits, AbsMean symmetric, and measure PPL.
\textbf{Maths.} If $F_L(x)=x+\Delta(x)$ with $\|\Delta\|$ small, the
layer's output is insensitive to bounded weight perturbations
$\|W-W_q\|_\infty\le\gamma$: the induced change in output is
$\le\|\Delta W\|\,\|x\|$, tiny next to $\|x\|$.
\textbf{Result.} SmolLM L5: 8-bit 17.76, 4-bit 17.77, 2-bit 17.76.
SmolLM L2: 19.11/19.07/19.18. GPT-2 flat at 38.25 (eval noise; larger
windows needed).
\textbf{Rule.} BI answers both questions: "can I delete this?" and "can I
crush its precision?".

\section{What failed, with the exact failure mode}

\subsection{02 - Affine bridge across the deleted gap}
We fit one ridge affine map $Y\approx AX+b$ from the input of the first
removed block to the output of the last and inserted it as a single
\texttt{nn.Linear}. \textbf{Why it fails.} The true map
$F_{last}\circ\dots\circ F_{first}$ is not affine (it is a composition of
LayerNorms, matmuls, and nonlinearities), so the single line cannot
approximate it; PPL jumped to $4.5\times 10^3$ (GPT-2 k=1) and
$1.9\times 10^8$ (SmolLM). \textbf{This is the direct motivation for
01.}

\subsection{08 - Parameter soup of adjacent blocks}
$\theta_{merge}=\alpha\theta_L+(1-\alpha)\theta_{L+1}$.\textbf{ Why it
fails.} Weight-space averaging across a neuron-permutation boundary is
meaningless; GPT-2 k=1 gave 103-252 PPL vs baseline 40.

\subsection{14 - Cross-layer weight tying}
Assigning \texttt{blocks[l+d] = blocks[l]} shares one set of weights
between two positions. \textbf{Why it fails.} Depth is not redundant in
weight space; tying at any period $d$ exploded PPL to $10^3$-$10^6\times$
baseline.

\subsection{04 - Permutation-aligned averaging}
Capture per-neuron activations of both adjacent blocks, run Hungarian
matching to align neurons, then average. \textbf{Why it fell back.} The
Hungarian step was not hardened up and the fallback was the naive average,
so the result equals 08 and inherits its failure.

\subsection{05, 10, 11, 12, 13 - scaffolds}
Each of LayerDrop training (11), learned L0 gates (10), attn/FFN pruning
proxy (12), early exit (13), and the 04-stage pipeline (05) executed but
produced flat/identical rows because the key training loop (exit-head
training, gate-sparsity thresholding, real LayerDrop fine-tune) was not
run inside the free MPS budget. They are kept in \texttt{archived/} as
incomplete.

\subsection{03, 09 - distillation at a tiny budget}
We wired KL-to-teacher (03) and KD-to-shallow-student (09) correctly but
the 60-step AdamW pass on MPS made the KL loss essentially zero, so
post-repair PPL equals pre-repair. The code path is reusable; only the
budget was missing.

\section{The working four compared head-to-head}
\small
\resizebox{\textwidth}{!}{\begin{tabular}{lllll}
\toprule
Approach & PPL at operating point & Params effect & Training? & Note \\
\midrule
01 Fitted-ternary TerGRASP & GPT-2 k=1 45.11, SmolLM 21.71 & $\sim$20$\times$ smaller on replaced block & yes, 80 steps & headline \\
06 BI baseline (targeted) & GPT-2 k=2 47.95, SmolLM 26.13 & block removed & none & control \\
07 SVD r=128 (SmolLM) & 22.00 vs 19.93 baseline & 2-4$\times$ on shrink blocks & none & smooth trade \\
15 2-bit on lowest-BI (SmolLM) & 17.76 vs 16.64 baseline & $\sim$16$\times$ on that block & none & rule \\
\bottomrule
\end{tabular}}
\normalsize

\section{Why 01 is the headline}
The other three working approaches each do one half of the job: 06 deletes
(no compression of the kept graph beyond removal), 07 shrinks in width but
keeps depth, 15 quantizes but keeps depth. 01 is the only one that
simultaneously (a) keeps the graph depth, (b) shrinks the replaced block
to $\sim$1.58 bit/param, and (c) holds PPL at the baseline operating point
$k{=}1$ on both models. Its failure mode is also the most defensible: at
$k\ge 3$ it tracks hard removal instead of exploding, so the claim is
bounded rather than brittle. That, plus the clean falsification control of
06 and the supporting rules of 07 and 15, is why we recommend 01.

\end{document}
